\runninghead{Stewardship Between Two Advents}
\properlane{vigilant-stewardship-between-two-advents}{Vigilant Stewardship of the Household Between the Two Advents}

A third interpretation begins with Mark~13:33--37 and asks how Jesus' command to watch orders the Church's life between Christ's two comings. Three concrete questions rise directly from his brief parable: why does the departing householder charge a single \emph{janitor} (``porter'' or doorkeeper) to watch while Jesus finishes by declaring, ``And what I say to you, I say to all: Watch''? Why does he name four watches of the Roman night --- evening, midnight, cockcrow, and morning? And how can a command to stay awake for the master's sudden arrival bind Christians in every century when only one generation will be alive at the Last Day? Five Fathers of the Church supply distinct parts of the answer: St.~Bede distinguishes clerical and lay watchfulness in Mark~13:34--37 (drawing on St.~Gregory the Great's Ascension and vigilance homilies), while St.~Gregory in Homily~13 (on Luke~12:35--40) maps the night watches onto the ages of human life; St.~Augustine (\work{Epistula}~199 §3, mediated through St.~Thomas Aquinas's \work{Catena Aurea in Marcum}) explains how the final day reaches every generation at death, and on Psalm~25(24):10 identifies the Lord's ways with Christ's two Advents; St.~Cyril of Jerusalem (\work{Catechesis}~15) contrasts Christ's first coming in humility with his second in glory; and St.~Leo the Great (\work{Sermon}~19) sets forth the Roman Advent discipline of temperance, fasting of the mind, and almsgiving.

\subsection{Office, Age, and the Hour of Departure in Mark~13}

St.~Bede (\work{In Marci Evangelium Expositio} IV.13) reads the household of Mark~13:34 as the ordered life of the Church after Christ's bodily Ascension, drawing his Ascension image from St.~Gregory the Great's Homily~9 (on Matthew~25:14--30) and his definition of watchfulness from Gregory's Homily~13 (on Luke~12:35--40). For Bede, the \emph{janitor} designates pastors and spiritual rulers charged with diligent care over the Church entrusted to them. Yet verse~37 (\emph{omnibus dico}, ``I say to all'') allows no layperson to delegate vigilance to the clergy: every Christian is commanded to guard the door of his own heart against evil suggestions. Quoting Gregory's Homily~13, Bede defines watching by three habits: keeping the eyes of the mind open to the true light, practicing in deeds what one believes, and driving off the darkness of sloth and negligence.

Gregory's Homily~13 also maps the three night watches of Luke~12:38 onto the seasons of a human life --- boyhood, youth, and old age --- showing that God is a \emph{patiens redditor} (``patient judge'') who leaves no age without hope of conversion yet requires stricter account if his patience is wasted, keeping the final hour hidden \emph{ut semper possit esse suspecta} (``so that it may always be anticipated''). In St.~Thomas Aquinas's \work{Catena Aurea in Marcum} (13:32--37), Theophylact of Ohrid distributes all four Roman watches of Mark~13:35 across life's arc (evening for old age, midnight for mid-youth, cockcrow for mature reason, morning for childhood) and reminds readers that Christ joins watching with prayer because wicked men also stay awake at night. In the same \work{Catena Aurea} passage, Pseudo-Jerome pictures sleep as chasing dream-phantoms that vanish upon waking --- echoing Gregory's warning in Homily~1 (on Luke~21:25--33, §§1, 5) not to set one's heart on a world visibly growing old.

St.~Augustine (\work{Epistula}~199 §3, via the \work{Catena Aurea in Marcum}) resolves the generational objection directly: because each believer's day of departure fixes the state in which he will be judged on the Last Day, personal death and the universal Parousia meet at every Christian's door, so no century can treat the Lord's return as someone else's concern.

\subsection{From Drowsiness to Watchful Service}

At the opening of the Mass, the psalmist of the Entrance Antiphon \emph{Ad te levavi} (Ps.~25[24]:1--3) tears his soul free from earthly ``briers and thorns'' (Bellarmine, Ps.~24, explanation~1) and carnal weight (Augustine, Ps.~24 §§1--2). Later in the same psalm (v.~10), Augustine identifies ``all the ways of the Lord'' --- mercy and truth --- with the two Advents of the Son of God, warning those freed by mercy in the first to keep watch before the severity of trial in the second. With the \work{Gloria} omitted (GIRM~53), the Collect (\emph{Da, quaesumus, omnipotens Deus}; modern text withheld) opens Advent prayer; its ancient Gelasian antecedent (II.lxxxiv) asks that God's family may go forth fittingly in righteous deeds (\emph{in operibus iustis aptos occurrere}) to meet the coming Christ --- words recalling Isaiah~64:5~Vg.\ (Heb.~64:4; \emph{Occurristi laetanti, et facienti justitiam}, ``Thou hast met him that rejoiceth, and doth justice'').

Against that active readiness, Isaiah~63:16b--64:8 (Vg.) portrays the spiritual drowsiness Mark warns against: a people drifting from God's ways (63:17), blown away like fallen leaves (64:6~Vg.), with ``none that riseth up, and taketh hold of thee'' (64:7~Vg.), even while remembering that God meets ``him that rejoiceth, and doth justice'' (64:5~Vg.). In the Responsorial Psalm (Ps.~80[79]), the assembly pleads for the Shepherd to stir up his might (v.~3b) and ``visit this vineyard'' (v.~15) --- a verb Bellarmine (Ps.~79, explanation~14) links to Jerusalem's tragedy in Luke~19:44 for failing to recognize ``the time of thy visitation'' --- and ends with a pledge of steady loyalty: ``And we depart not from thee... we will call upon thy name'' (v.~19). In the Second Reading (1~Cor.~1:3--9), Paul assures the Corinthians that they lack no spiritual gift while awaiting Christ's manifestation (1:7), yet as St.~John Chrysostom observes (Homily~II), promising to confirm them ``unto the end without crime'' (1:8) quietly warns a church threatened by party quarrels; read alongside Mark~13:34, Christ's household is already equipped for its work yet cannot afford internal division. Before the Gospel, the Gospel Acclamation \emph{Ostende nobis, Domine} (Ps.~85[84]:8) asks the Lord to show his mercy and grant his salvation.

\subsection{Eucharistic Consecration and the Guard Against Moralism}

In the Liturgy of the Eucharist, daily stewardship is consecrated to the final kingdom. In the Prayer over the Offerings (\emph{Suscipe, quaesumus, Domine}; modern text withheld), the ancient Veronense antecedent (XXVIIII.372) offers temporal devotion (\emph{devotioni... temporalis}) so that it may become the reward of eternal redemption. At the Seasonal Preface of Advent (GIRM~364; modern text withheld), Eucharistic thanksgiving stands inside the two-Advent horizon expounded by St.~Cyril of Jerusalem (\work{Catechesis}~15 §§1--2, 4, 26): contrasting Christ's first coming in patience, swaddling clothes, the manger, and the Cross with his second in light, angels, glory, and the crown of a divine kingdom (citing Malachi~3:1--2), Cyril forbids curiosity about dates (Acts~1:7) and reduces readiness to plain works of mercy from Matthew~25:35--41 (``there is here no need of allegory, but to fulfil what is said''), keeping one's lamp constantly burning. At Holy Communion, the Communion Antiphon \emph{Dominus dabit benignitatem, et terra nostra dabit fructum suum} (Ps.~85[84]:13) speaks of the Lord giving goodness so that our earth yields fruit --- a fruit St.~Augustine describes in sobriety, giving, prayer, and hymns (Ps.~84 §12), matching St.~Leo the Great's Roman Advent discipline (\work{Sermon}~19 §§1--3, citing Luke~21:34) of pairing bodily temperance with fasting of the mind from vice and generous almsgiving to the poor. The Prayer after Communion (\emph{Prosint nobis}; modern text withheld) closes the Communion Rite, and the Concluding Blessing (GIRM~90) sends the assembly out at the dismissal to live in good works and praise of God.

Putting Mark~13:33--37 and active stewardship in the foreground brings one clear pastoral hazard: if detached from the confession of frailty in Isaiah~64:6--8~Vg.\ (Heb.~64:5--7) and the priority of grace in 1~Corinthians~1:4--9, an exhortation to keep busy in good works can drift toward Pelagian self-effort or neglect the consolation needed by parishioners bewildered by God's apparent silence.

\subsection{The Four Senses of Scripture}

\begin{fourSenses}
\item[Literal.] Isaiah laments a complacent, drifting people while recalling that God meets those who do justice (Isa.~63--64); the psalmist begs God to visit his vineyard and binds the people not to turn away (Ps.~80[79]); Paul urges a gifted urban congregation to stand blameless until the Day of Christ (1~Cor.~1:3--9); Jesus on the Mount of Olives delivers his parting charge to keep watch (Mark~13:33--37).
\item[Allegorical.] The household committed to servants and a doorkeeper (Mark~13:34, Bede) and the vineyard visited by God (Ps.~80[79]:15, Bellarmine and Augustine) signify the pilgrim Church structured with pastors and faithful living between Christ's First Advent in humility and his Second Advent in glory (Cyril of Jerusalem, \work{Catechesis}~15 §§1--2; Augustine, \work{Enarr.}, Ps.~24 §10).
\item[Moral.] Watching means concrete daily fidelity --- keeping the mind open to truth, practicing what one believes, driving off torpor (Gregory, Hom.~13, quoted by Bede), and uniting bodily temperance with fasting from vice and almsgiving to the poor (Leo, \work{Sermon}~19) across every stage of life from childhood to old age (Gregory, Hom.~13 on Luke~12:38; Theophylact on Mark~13:35).
\item[Anagogical.] Each believer's hour of death (Augustine, \work{Ep.}~199 §3 via Aquinas's \work{Catena Aurea}) and the universal Parousia (1~Cor.~1:8; Cyril, \work{Catechesis}~15 §§1, 26) meet in the sudden return of the Master, when servants found awake in works of justice are gathered to Christ's right hand (\work{Gelasianum Vetus} II.lxxxiv antecedent) to inherit the heavenly kingdom.
\end{fourSenses}
